% We introduce the \framework framework, a self-referential agent capable of self-awareness and dynamic self-modification of its logical code, guided by both environmental feedback and decisions generated by LLMs. 
% This architecture facilitates recursive self-improvement, allowing the agent to iteratively optimize its performance. 
% Inspired by the self-referential capabilities of the \godel machine, the \framework adapts this concept to modern LLMs and contemporary agent design paradigms. 
% A key distinction from the original \godel machine is the replacement of the proof search mechanism with an LLM, which provides generalized reasoning and decision-making abilities. 
% Moreover, rather than relying on expected utility as the primary criterion for code modification, the agent's decisions are driven by the flexibility and reasoning capacity of the LLM. 
% This allows the agent to selectively determine when environmental feedback is necessary, rather than requiring constant interaction with the environment at each step.

\begin{algorithm*}[t!]
\small
\caption{Recursive Self-Improvement of \framework}
\vspace{-12pt}
\label{alg:1}
\begin{multicols}{2}
\begin{algorithmic}[1]
\State \textbf{Input:} Initial agent policy $\pi_0$, initial decision function $f_0$, goal $g$, environment state $\mathcal{E}$, utility function $U$, self code reading function \texttt{SELF\_INSPECT}
\State \textbf{Output:} Optimized policy $\pi$ and \framework $s$
\State \Comment{Get all agent code, including the code in this algorithm.}
\State $s \gets \texttt{SELF\_INSPECT}()$
\State \Comment{Compute the initial performance.}
\State $r \gets U(\mathcal{E}, \pi_0)$
\State \Comment{Perform recursive self-improvement.}
\State $\pi, s \gets \texttt{SELF\_IMPROVE}(\pi, s, r, g)$
\State \Return $\pi, s$
\State \Comment{Initial code of self-referential learning.}
\Function{\texttt{SELF\_IMPROVE}}{$\mathcal{E}, \pi, s, r, g$}
    \State \Comment{Obtain action sequence.}
    \State $a_1, \ldots, a_n \gets f_0(\pi, s, r, g)$
    \For{$a_i$ \textbf{in} $a_1, \ldots, a_n$}
        \State $\pi, s, r \gets \texttt{EXECUTE}(\mathcal{E}, \pi, s, r, a_i)$
    \EndFor
    \State \Return $\pi, s$
\EndFunction
\State \Comment{Initial action execution function.}
\Function{\texttt{EXECUTE}}{$\mathcal{E}, \pi, s, r, a$}
    \State \textbf{switch} $a.\texttt{name}$
    \State \quad \textbf{case} \texttt{self\_state}:
    \State \quad\quad $s \gets \texttt{SELF\_INSPECT}()$
    % \quad \Comment{Update self-awareness information of the policy.}
    \State \quad \textbf{case} \texttt{interact}:
    \State \quad\quad $r \gets U(\mathcal{E}, \pi)$ 
    % \quad \Comment{Update performance based on agent-environment interaction.}
    \State \quad \textbf{case} \texttt{self\_update}:
    \State \quad\quad $\pi, s \gets a.\texttt{code}$ 
    % \quad \Comment{Modify policy logic based on the action.}
    \State \quad \textbf{case} \texttt{continue\_improve}:
    \State \quad\quad \Comment{Recursively invoke self-improvement.}
    \State \quad\quad $\pi, s \gets \texttt{SELF\_IMPROVE}(\mathcal{E}, \pi, s, r, g)$ 
    \State \Return $\pi, s, r$
\EndFunction
\vspace{-10pt}
\end{algorithmic}
\end{multicols}
\end{algorithm*}





In this section, we first describe the formal definitions for previous agent methods with a lower degree of freedom, including hand-design and meta-learning optimized agents, as a background. Then we introduce our proposed \framework, a self-referential agent that can recursively update its own code, evolving over training.

% In this section, we describe our proposed \framework, a self-referential agent that can recursively update its own code. Firstly, we describe normal agents with a lower degree of freedom, which have a fixed algorithm, as a background. Then we describe our \framework, whose algorithm is evolving over training. We describe a possible initial algorithm of \framework in \cref{alg:1}.

% We first provide formal definitions for hand-design agents and meta-learning optimized agents as a background to facilitate the understanding of \framework.

Let $\mathcal{E} \in \mathcal{S}$ denote a specific environment state, where $\mathcal{S}$ denotes the set of all possible environments the agent will encounter. For example, an environment can be a mathematical problem with ground truth solutions.
%\wxy{not sure if this is accurate. double check.}.
We denote the policy that an agent follows to solve a problem in the current environment by $\pi \in \Pi$, where $\Pi$ is the set of all possible policies the agent can follow. 

% paragraph
A \textbf{hand-designed agent}, as shown in the left panel of Figure \ref{fig:compare}, is not capable of updating its policy and following the same policy $\pi$ all the time, regardless of environmental feedback.

In contrast, a \textbf{meta-learning optimized agent} updates its policy based on a meta-learning algorithm $I$ at training time based on the feedback it receives from the environment, as shown in the middle panel of Figure \ref{fig:compare}. The environment feedback is usually defined as a utility function $U: \mathcal{S} \times \Pi \rightarrow \mathbb{R}$, which maps an environment and a policy to a real-valued performance score. The main training algorithm of a meta-learning optimized agent can then be written as follows:
\begin{align*}
    \pi_{t+1} = I(\pi_t, r_t),\;\;\; r_t = U(\mathcal{E}, \pi_t),
\end{align*}
In this case, the agent's policy $\pi_t$ evolves at training time, with the learning algorithm $I$ updating the policy based on feedback $r_t$, while the meta-learning algorithm $I$ remains fixed all the time.



A \textbf{self-referential \framework}, on the other hand, updates both the policy $\pi$ and the meta-learning algorithm $I$ recursively. The main idea is that, after each update, the whole code base of the agent is rewritten to accommodate any possible changes. Here we call this self-updatable meta-learning algorithm $I$ a self-referential learning algorithm. The training process of a \framework can then be written as:
%Specifically, the process is defined by the following function: $\pi_{t+1} = \pi_t(\mathcal{E}, \pi_t, g)$, where $g \in \mathcal{G}$ represents the goal task.
% \begin{align}
%     \pi: \mathcal{E} \times \Pi \times \mathcal{G} \rightarrow \Pi,
% \end{align}
% where $\mathcal{G}$ represents the set of goals.
\begin{align*}
    \pi_{t+1}, \; I_{t+1} = I_t(\pi_t, I_t, r_t, g),\;\;\; r_t = U(\mathcal{E}, \pi_t),
\end{align*}
where $g \in \mathcal{G}$ represents the high-level goal of optimization, for example, solving the given mathematical problem with the highest accuracy. Such a recursive design of the agent requires the specification of an initial agent algorithm $(\pi_0, I_0)$, detailed as follows:
%The goal can be general, such as "maximize utility within the given environment."
% To facilitate recursive self-improvement and integrate with modern agent architectures, the \framework adopts a decision-making and action paradigm. The agent's policy set, $\Pi = \{(f, o) \mid f \in F, o \in O\}$, is composed of two key functions, which define each policy $\pi \in \Pi$:
\begin{itemize}[leftmargin=*,itemsep=1pt]
    \item A initial agent policy $\pi_0$ to perform the desired task within the environment $\mathcal{E}$. For example, it can be chain-of-thought prompting of an LLM.
    \item A self-referential learning algorithm $I_0$ for recursively querying an LLM to rewrite its own code based on the environmental feedback.
\end{itemize}
We then further specify a possible initialization of the self-referential learning algorithm $I_0 = (f_0, o_0)$, using a mutual recursion between a decision-making function $f_0$, and an action function $o_0$:
\begin{itemize}[leftmargin=*,itemsep=1pt]
    \item The decision-making function $f_0$, implemented by an LLM, determines a sequence of appropriate actions $a_1, a_2, ..., a_n \in \mathcal{A}$ based on the current environment $\mathcal{E}$, the agent's algorithm $(\pi_t, I_t)$, and the goal $g$.
    \item The action function $o_0$, executes the selected action and updates the agent's policy accordingly. 
    % If no suitable action can be executed, call the decision-making function $f$ again for new actions.
\end{itemize}



The set of actions $\mathcal{A}$ for the action function $o$ to execute needs to include the following four actions: 
\begin{itemize}[leftmargin=*,itemsep=1pt]
    \item \texttt{self\_inspect}: Introspect and read the agent's current algorithm $(\pi_t, I_t)$.
    \item \texttt{interact}: Interact with the environment by calling the utility function $U$ to assess the performance of the current policy $\pi_t$. 
    \item \texttt{self\_update}: Alter and update $(\pi_t, I_t)$ with an LLM and produce $(\pi_{t+1}, I_{t+1})$.
    \item \texttt{continue\_improve}: If no other actions can be taken, recursively invoke the decision algorithm $f$ to produce new actions.
\end{itemize}
The agent code is updated to $(\pi_{t+1}, I_{t+1})$ after the current execution of $(\pi_t, I_t)$ is finished.
Both the agent algorithm $(\pi, I)$ and the action set $\mathcal{A}$ are not static and can be expanded and modified by the agent itself at the training time.
Algorithm \ref{alg:1} illustrates the described algorithm for the \framework. Each recursive call enables the agent to refine its logic and become progressively more efficient.

